% English version of paper/ja_v2/02_関連研究.md (2026-11-17).
\section{Related work}\label{sec:related}

\subsection{Semi-supervised domain adaptation}
Methods usable in semi-supervised domain adaptation, in which the target has a few labels, have been proposed for each information source: learning representations that do not reveal the domain (domain-adversarial neural networks, DANN) \cite{ganin2016dann} and aligning feature covariances (CORAL) \cite{sun2016coral} for unlabelled targets, entropy minimisation and pseudo-labels from semi-supervised learning \cite{grandvalet2004,lee2013pseudolabel}, and minimax entropy (MME) for semi-supervised domain adaptation \cite{saito2019mme}. Adding source data can also lower target performance (negative transfer) \cite{wang2019negative}. The method studies above compare methods that combine information sources and do not measure, in one setting, which part of the classifier error each source corrects.

\subsection{Cross-corpus MER with few labels}
For audio emotion classifiers that degrade on another corpus, adversarial learning that aligns representations with unlabelled target utterances \cite{abdelwahab2018dann} and learning that pulls together representations learned per corpus \cite{gideon2021addog} have been proposed. Ryumina et al.\ \cite{ryumina2025} evaluated the four corpora of our leave-one-corpus-out setting (IEMOCAP, MOSEI, MELD and AFEW; we use only the improvised part of IEMOCAP) with four emotions without target labels and showed that training on several corpora raised UAR by 1.7--6.6\pt. Lian et al.\ \cite{lian2026merbench} built cross-corpus benchmarks with shallow classifiers on frozen features and reported that multimodal models exceeded unimodal ones by small differences across corpora. For targets with a few labels, Feng and Chaspari \cite{feng2023} proposed metric learning for audio features, Li et al.\ \cite{li2021acted} adaptation of audio and video between acted and improvised speech, and He et al.\ \cite{he2025m4ser} a method using audio and speech-recognition transcripts.

\subsection{Distillation from other modalities}
Albanie et al.\ \cite{albanie2018} attached facial-expression predictions to unlabelled videos and trained an audio emotion classifier on them, and Zhang et al.\ \cite{zhang2021combining} combined this scheme with semi-supervised learning. Transferring representations or predictions of one modality to a model of another through paired unlabelled data is also used outside emotion recognition \cite{gupta2016,jaritz2020xmuda}. Training on the predictions of a more accurate model of another modality has been reported not to always improve the trained model \cite{xue2023mfh}.

\subsection{Late fusion and reliability weights}\label{sec:rel-fusion}
Zadeh et al.\ \cite{zadeh2017tfn} criticised late fusion, which trains unimodal classifiers separately and combines their predictions, for not learning interactions between modalities because it assumes that simple weighted averaging is a proper fusion approach. Predictions can be combined by product or sum rules \cite{kittler1998}. Busso et al.\ \cite{busso2004} reported that modality strengths differ by emotion, Kaya et al.\ \cite{kaya2017} chose model-by-emotion weights on AFEW by random search on validation data, M3ER \cite{mittal2020m3er} learned a multiplicative fusion layer that strengthens reliable modality cues per utterance, and Mar\'in et al.\ \cite{marin2026classwise} used per-utterance, per-emotion modality weights output by a network trained with all labels within one corpus.

\subsection{Position of this study}
We measure, with one decomposition, which part of the error each of the four sources corrects (RQ1(a)), and then examine on four corpora which classifiers distillation improves and when in cross-corpus MER with a few labels per emotion (RQ1(b)). Unlike Albanie et al., who used facial predictions only, and Zhang et al., who mixed facial and audio predictions, our pseudo-labels are built from three classifiers including text. In RQ2, we use late fusion (Sect.~\ref{sec:rel-fusion}) to combine predictions into pseudo-labels and measure one weight per classifier and emotion. We chose decision-level fusion because, in our analysis, the strongest modality differed by corpus and emotion, and classifier-by-emotion weights can represent this difference using a few labels and the source corpora; multimodal models trained on few samples can also fall below the best unimodal model \cite{huang2021}. The proposal measures strengths on the K labels and per-emotion profiles on the source corpora and, unlike M3ER and Mar\'in et al., uses one weight per classifier and emotion. To our knowledge, splitting a reliability table into strengths measured on the K target labels and profiles measured on the source corpora has not been reported.
